% Benchmark scope and verified local dataset contracts.
% Compile through main.tex.

\section{The benchmark landscape for one-to-many simulation}
\label{sec:bench}

\subsection{Match the dataset to the statistical claim}
\label{sec:bench:criterion}

A distributional task must specify the model-visible condition $g$, the source
of random variation and the outcome being scored. Repeated outcomes at fixed
$g$ permit direct empirical checks of that conditional distribution. They are
not a prerequisite for proper-score evaluation: one outcome at each of many
test conditions can estimate an average predictive score. Such an evaluation
does not, without additional assumptions or observations, resolve calibration
separately at every individual condition~\cite{gneiting2007strictly}.

We therefore distinguish the suitability of a dataset for a specific
hidden-variable experiment from its validity as a probabilistic benchmark.
Table~\ref{tab:landscape} identifies the different claims supported by common
sampling designs. The examples below are a targeted review and an audit of
three local conversions, not an exhaustive census of all available datasets.

\begin{table}[t]
\centering
\small
\caption{Sampling designs and their interpretation. Representation, conditioning
and dependence must be specified separately from the number of samples.}
\label{tab:landscape}
\begin{tabular}{@{}p{0.27\linewidth}p{0.32\linewidth}p{0.32\linewidth}@{}}
\toprule
\textbf{Design} & \textbf{Can evaluate} & \textbf{Main limitation} \\
\midrule
One outcome per condition &
Average proper score across test cases &
Few direct checks of each individual conditional law \\
Replicates at many conditions &
Within-condition distributions and transfer between conditions &
Independent reference samples and a declared condition split are needed \\
Many outcomes at one condition &
A distribution at that fixed condition &
No evidence of transfer between conditions \\
Snapshots from one long run &
A stationary distribution under justified sampling assumptions &
Serial dependence reduces effective sample size \\
Several solver or model variants &
A declared mixture over model choices &
Not automatically physical replicate variability \\
\bottomrule
\end{tabular}
\end{table}

Mesh-based generative fluid benchmarks already contain multiple physical
states for varying geometries~\cite{lino2025diffusiongraphnets,
lino2026scaleautoregressive}. Those states represent statistically developed
flow and need not be independent nonlinear solves with a withheld imperfection.
The distinction motivates our shell construction; it does not justify claiming
that no many-condition, many-outcome mesh benchmark exists.

\subsection{Examples from a public simulation collection}
\label{sec:bench:tier1}

The Well includes several useful seed ensembles, with different conditioning
contracts~\cite{ohana2024well}. Its source documentation must be read at the
level of each subset rather than assuming that every trajectory in the
collection has the same uncertainty mechanism.

\paragraph{Turbulence--gravity--cooling.}
The documented ensemble has $2{,}700$ trajectories: $100$ random seeds for
each of $27$ combinations of initial density, temperature and metallicity.
The public fields are supplied on a structured grid. These counts
describe the simulation design; an evaluation must still declare which initial
fields are exposed to the model and align physical output
times.\footnote{\url{https://polymathic-ai.org/the_well/datasets/turbulence_gravity_cooling/}}

\paragraph{Rayleigh--B\'enard convection.}
The documented design combines five Rayleigh numbers, seven Prandtl numbers,
five initial buoyancy amplitudes and ten noise seeds, giving $1{,}750$
trajectories. If all three non-seed parameters are visible, there are $175$
conditions with ten seeds each. Grouping only by Rayleigh and Prandtl number
instead gives $35$ groups and marginalises over the five buoyancy amplitudes;
these are different prediction tasks. The collection also distinguishes the
original graded grid from a uniformly resampled
version.\footnote{\url{https://polymathic-ai.org/the_well/datasets/rayleigh_benard/}}

\paragraph{Turbulent radiative layer.}
The two-dimensional subset supplies ten seeds for each of nine cooling times.
Our local conversion retains $72$ training trajectories and $18$ inference
trajectories. Its state channels include the velocity perturbation that varies
between seeds; Section~\ref{sec:bench:negative} explains why the visible input
changes the interpretation of this ensemble.

\paragraph{Viscoelastic instability.}
The source describes coexisting dynamical states. A multimodal field law
motivates richer models than a single Gaussian distribution over the output
field, but does not prove that flow matching is necessary: nonlinear latent
decoders and other mixture or generative models can also represent multiple
modes. Any comparison must specify the distribution of initialisations and how
trajectory lengths affect mode
weights.\footnote{\url{https://polymathic-ai.org/the_well/datasets/viscoelastic_instability/}}

None of these counts gives a universal threshold for reliable CRPS or rank
histograms. Precision depends on the target distribution, effect size,
dependence and aggregation rule, as well as the number of generated members
and reference outcomes.

\subsection{Common mismatches between data and claims}
\label{sec:bench:reject}

Several distinctions prevent misleading benchmark conclusions.

\begin{itemize}
  \item \textbf{Visible variability versus hidden variability.}
    Two runs with different material maps or initial velocity fields do not
    share a model-visible condition when those fields are inputs. They can
    still form a valid conditional prediction dataset, but are not replicate
    outcomes at identical inputs.
  \item \textbf{Time samples versus independent realisations.}
    Correlated snapshots may estimate stationary statistics under appropriate
    stationarity and ergodicity assumptions. Treating them as independent
    replicates overstates precision; it does not make the stationary
    distribution itself invalid.
  \item \textbf{An ensemble mean versus a realised field.}
    If the released target is an average over solver runs, its distribution is
    the distribution of that average. Evaluating it does not recover the
    variability of individual runs.
  \item \textbf{Physical variability versus model uncertainty.}
    Alternative constitutive models, solver settings and neural training seeds
    can define useful uncertainty ensembles. Their mixture and weighting must
    be explicit, and should not be labelled independent physical repeats.
\end{itemize}

Mechanical MNIST Crack Path illustrates the first distinction: the original
problem relates microstructural input to mechanically driven
fields~\cite{mohammadzadeh2022crackpath}. Withholding that input changes the
task, but does not automatically make the remaining observed state identical
between cases.

\subsection{A withholding construction and its input boundary}
\label{sec:bench:construction}

The local Mechanical MNIST conversion omits the inclusion field and stores
coordinates, damage and two displacement components on a $256\times256$
regular grid, connected by non-periodic grid edges. It contains $1{,}750$
training cases and $250$ inference cases, each with $20$ stored frames. These
are grid samples of finite-element results, not the original finite-element
mesh. The construction differs from the shell dataset of
Section~\ref{sec:results:buckling}, where the visible and hidden imperfection
components are specified before simulation.

Crucially, the converted crack dataset does \emph{not} have bit-identical
model-visible initial states. A direct comparison of the first two numerically
indexed cases in each HDF5 split finds differences in all three state channels
at the first stored frame. The maximum absolute damage difference is
approximately $0.00881$ in that training pair and $0.0780$ in that inference
pair; both displacement channels differ as well. These counterexamples are
enough to reject identical-input replication, but do not measure how much
information about the final crack is available from that frame.

Thus, $p(y_{\mathrm{future}}\mid \text{stored state})$ and
$p(y_{\mathrm{trajectory}}\mid \text{geometry, loading protocol})$ are
different targets. Omitting the explicit inclusion channel does not remove the
inclusions' mechanical effects from the observed state. Under the current
autoregressive input contract, the stored state is part of the condition.
A genuinely common-state construction would require defining and validating a
common initial input, together with a compatible loading-time convention.
That variant is not supplied by simply dropping the inclusion channel.

The local audit establishes this input-contract limitation. It does not
establish a numerical information-leak floor, the residual conditional
dimension, or a precise branching window. Such results require a reproducible
analysis with specified sample selection, split membership, distance
normalisation and uncertainty estimates. Per-step diagnostics would be
appropriate for locating any branching interval, but are not presented here
as completed measurements.

\subsection{Two candidate slots with different targets}
\label{sec:bench:negative}

\paragraph{A radiative-layer seed ensemble.}
The local radiative-layer files contain nodal arrays of shape
$[8,101,49152]$: three coordinates, density, pressure, two velocity components
and cooling time. Each of the nine cooling times has eight training seeds and
two inference seeds. Comparing two seeds at the same cooling time shows
bit-identical initial density and pressure but different initial velocities.
The model therefore receives at least part of the seed-dependent perturbation.

For a deterministic solver with a complete state, conditioning on that state
removes uncertainty about the next step. This observation does not prove that
the stored fields are a complete Markov state or that a learned predictor has
zero residual uncertainty. It does show why a cross-seed ensemble conditioned
only on cooling time cannot be used directly as the reference conditional law
for a model that also observes the realised velocity field.

The two reference seeds per cooling time are a small sample for estimating a
condition-specific distribution. This reference count does not determine the
number of rank-histogram bins: a histogram for an ensemble of $M$ generated
members has $M+1$ possible ranks, with tie handling specified. Reference count
controls how many rank observations are available, and nodes or timesteps
within a trajectory cannot simply be counted as independent observations.

\paragraph{Grain growth and the meaning of a Brier score.}
The local SPPARKS conversion has $900$ training and $100$ inference samples,
each storing a static $100^3$ grain-identifier array. Numeric grain IDs are
arbitrary labels: a bijective relabelling leaves the partition unchanged
while altering a squared error on the IDs. A separate crop converter provides
a $50^3$ boundary-indicator representation using open faces and a
six-neighbour stencil. The invariant target resolves label symmetry, although
it does not alone specify an adequate morphology metric.

A spatially constant \emph{probability forecast} must not be confused with a
constant generated microstructure. If a boundary event at a node has
probability $p$, the expected Brier score of a forecast probability $q$ is
\begin{equation}
  \mathbb{E}(q-Y)^2 = p(1-p)+(q-p)^2,\qquad Y\sim\operatorname{Bernoulli}(p).
  \label{eq:brier-marginal}
\end{equation}
It is minimised at $q=p$. If a probability is instead estimated as the mean
$\hat q$ of $M$ independent Bernoulli samples with success probability $q$,
independent of the verifying outcome, the expected score gains the term
$q(1-q)/M$. A known probability can therefore outperform the empirical
probability from a finite ensemble sampled from the same law. This is the
finite-ensemble effect addressed by fair scores~\cite{ferro2014fair}, not
evidence against Brier propriety.

If two field distributions have identical probabilities at every node,
summing marginal Brier scores cannot distinguish their spatial dependence.
Scrambling can preserve marginal probabilities while destroying grains.
The appropriate conclusion is to supplement valid marginal probability
scores with joint or morphology-sensitive diagnostics, not to call marginal
scores anti-informative. A proper score applied to a descriptor assesses its
induced distribution; it is strictly informative about the full field only
when the representation and score identify that field law.

At fixed physical settings this is a legitimate distributional task. It can
measure variability at those settings, but supplies no test of conditioning
on varying geometries or process parameters.

\subsection{Design requirements for the shell benchmark}
\label{sec:bench:missing}

The shell construction targets a particular engineering question: how a
withheld imperfection changes the field distribution across known geometries.
Its requirements are (i)~a declared visible-input contract and independently
drawn hidden variables; (ii)~retained mesh connectivity and consistently
defined output fields; (iii)~repeated outcomes and separate condition splits
for within-condition evaluation and transfer; (iv)~documented solver,
convergence and retention criteria; and (v)~engineering quantities of interest
alongside field diagnostics. A deterministic control is useful for detecting
spurious spread, but a matching spread on that control alone cannot certify
calibration.

Sample counts should follow precision or power requirements rather than a
universal cutoff such as $50$ conditions with $50$ members each. Finite
reference self-scores provide a sampling baseline, not a strict lower bound
for every observed score. Likewise, exclusions based on numerical convergence
may change the retained conditional law and must be reported.
Section~\ref{sec:results:buckling} states the implemented construction and
the current evidence for these requirements.

